\chapter{项目概述}

\section{项目背景与问题分析}

在个人财务管理场景中，自动记账工具（如微信账单导出、支付宝账单导出）已成为日常财务整理的重要手段。然而，传统记账软件在自动分类环节存在显著短板，直接影响统计分析的准确性。一个典型且普遍的场景是：当用户向校园一卡通充值时，账单上的商户名称可能仅显示为「东南大学」。基于关键词匹配的传统分类引擎很容易将其归入「学习」类别，但从真实消费意图来看，这笔交易更接近于「餐饮/校园卡充值」。这种语义误判一旦发生，不仅会扭曲用户的消费结构认知，还会在没有校正机制的情况下持续污染后续的财务统计。

进一步分析，个人用户在整理微信、支付宝账单时通常会遇到三类共性问题：

\begin{enumerate}
  \item \textbf{自动分类依赖浅层规则}：传统分类引擎仅根据商户名或关键词做简单匹配，难以理解「在便利店买口罩」和「在便利店买零食」之间的消费意图差异。
  \item \textbf{分类依据不透明}：用户不知道一笔交易为什么被归入某个类别，缺乏解释性，因而无法判断分类是否可信。
  \item \textbf{误判缺少校正入口}：当分类出错时，用户要么被动接受错误结果，要么需要逐个手动修改，缺乏一个高效、可控的人工校正机制。
\end{enumerate}

Boring Financial 正是针对上述问题设计的解决方案。系统围绕“导入—大模型语义初分—人工校正—多维度分析—报表输出”的完整业务流程展开，将原始账单转换为语义准确、可解释、可修正的财务分析结果，同时引入消费人格画像和家庭组织聚合等分析维度，为用户提供比传统记账工具更深入的财务洞察。

\section{项目目标}

Boring Financial 的核心定位是一个面向个人与家庭的智能账单分类与财务分析系统。项目的具体目标如表~\ref{tab:project-goals} 所示。

\begin{table}[H]
\centering
\caption{项目目标}
\label{tab:project-goals}
\begin{tabular}{p{3.8cm}p{9.4cm}}
\toprule
\textbf{目标} & \textbf{说明} \\
\midrule
完成记账分析流程 & 支持账单导入、格式解析、大模型语义初分、人工校正、统计看板、消费人格画像、家庭组织聚合和 PDF 报表的完整流程。 \\
\addlinespace
支持多用户与数据隔离 & 用户默认只能访问自己的账单、交易、分类和报表；加入家庭组织后可按成员角色查看聚合数据，个人数据的独立性不受影响。 \\
\addlinespace
支持语义分类与人工校正 & 以大模型语义分类为主要手段，同时保留规则分类作为后备方案，并为低置信度交易提供人工校正入口，确保用户对最终分类结果拥有控制权。 \\
\addlinespace
支持工程化部署 & 提供清晰的系统架构、接口文档、开发指南、测试套件和部署说明，支持从本地开发到 Docker Compose 裸机部署的多种运行模式。 \\
\addlinespace
控制实现复杂度 & 聚焦记账分析业务流程本身，不引入支付、清算和银行级风控等超出项目范围的功能。 \\
\bottomrule
\end{tabular}
\end{table}

\section{项目范围}

\subsection{范围内功能}

\begin{table}[H]
\centering
\caption{项目范围内功能模块}
\label{tab:in-scope}
\begin{tabular}{p{3cm}p{10.2cm}}
\toprule
\textbf{模块} & \textbf{范围说明} \\
\midrule
用户认证 & 注册、登录、JSON Web Token（JWT）身份鉴别、多用户数据隔离。 \\
账单导入 & 支持微信、支付宝账单文件上传，管理导入批次状态和导入记录。 \\
账单解析 & 将不同平台的原始账单格式转换为统一的交易数据结构，支持去重。 \\
大模型初分 & 支持规则分类、缓存分类和外部大模型语义分类的组合处理流程，提供分类理由与置信度，并设置失败重试机制。 \\
人工校正 & 用户可筛选和复核低置信度交易，修正错误分类，修正结果在后续统计中体现。 \\
分类管理 & 支持系统预设分类和用户自定义分类，系统分类不可禁用。 \\
财务看板 & 展示收入、支出、结余、储蓄率、分类占比分布、Top 商户和收支趋势。 \\
消费人格 & 基于行为经济学理论生成 16 型消费人格画像、五维财务健康评分，并提供自评问卷与数据画像的偏差对比。 \\
家庭组织 & 支持创建家庭组织、管理所有者、管理员和普通成员角色，聚合家庭财务看板、家庭消费人格和家庭报表。 \\
PDF 报表 & 按日期范围和导入文件生成可下载的 PDF 财务报告。 \\
系统部署 & 支持本地开发运行、Docker Compose 多服务部署和裸机 Nginx 与 systemd 部署。 \\
\bottomrule
\end{tabular}
\end{table}

\subsection{范围外内容}

\begin{table}[H]
\centering
\caption{项目范围外内容及原因}
\label{tab:out-scope}
\begin{tabular}{p{4.2cm}p{9cm}}
\toprule
\textbf{不做内容} & \textbf{原因} \\
\midrule
支付、转账、清算 & 涉及真实资金流和金融合规要求，超出项目范围。 \\
银行级风控与审计 & 需要严格的合规流程和高安全保障，不属于记账分析业务流程的核心关注点。 \\
多机构财务管理 & 当前定位为个人与家庭的账单分类与分析工具。 \\
复杂家庭预算审批流程 & 当前家庭组织模块聚焦成员聚合和角色管理，不做多级审批流。 \\
商业化运营后台 & 系统聚焦个人使用和轻量部署，不覆盖商业运营能力。 \\
\bottomrule
\end{tabular}
\end{table}

\section{项目收益}

Boring Financial 的收益可以从多个利益相关方视角进行审视，如表~\ref{tab:benefits} 所示。

\begin{table}[H]
\centering
\caption{项目收益分析}
\label{tab:benefits}
\begin{tabular}{p{2.8cm}p{10.4cm}}
\toprule
\textbf{受益方} & \textbf{收益} \\
\midrule
普通用户 & 减少误分类带来的统计偏差，快速理解个人支出结构、消费趋势和消费行为特征。通过消费人格画像获得对自身消费心理的深层认知，通过家庭组织模块了解家庭成员的整体财务状况。 \\
\addlinespace
家庭管理员 & 可将个人账本聚合为家庭视图，了解家庭整体的收入支出情况和消费结构，无需为家庭成员创建共享账本或手动合并数据。 \\
\addlinespace
系统运维人员 & 基于清晰的模块边界、标准化的接口文档、自动化测试套件和多种部署方案，可快速搭建和维护系统运行环境。 \\
\addlinespace
项目开发团队 & 在实践中覆盖了需求分析、系统设计、编码实现、测试验证、部署运维和文档整理的全流程，积累了从单体应用到前后端分离系统的完整开发经验。 \\
\bottomrule
\end{tabular}
\end{table}
